% ============================================
% Section: System Model
% ============================================
\label{sec:system}

Co-LIC2 consists of $N$ \emph{nodes} cooperating over a wireless network. Each
node is a moving platform that runs a generalized Gaussian-LIC2 pipeline
locally and exchanges compressed submaps with peers. This section formalizes
the node model, the per-node multi-sensor fusion, the network topology, and the
overall architecture.

\subsection{Node Model}
\label{sec:node}

A node $i$ carries:
\begin{itemize}
  \item $N_c^{(i)} \geq 1$ cameras with partially overlapping fields of view;
  \item $N_l^{(i)} \geq 1$ LiDARs of possibly different types (mechanical
    spinning, solid-state, MEMS);
  \item one or more IMUs, fused into a single virtual body-frame IMU;
  \item an onboard edge GPU (e.g., Jetson Orin, RTX 4090);
  \item a wireless interface (Wi-Fi 6/6E, 5G NR, or ad-hoc mesh).
\end{itemize}

Each node defines a \emph{body frame} $B_i$ anchored at its IMU center. Every
sensor extrinsic is expressed as a rigid transform $T_{B_i \leftarrow S}$ from
the sensor frame $S$ to $B_i$. The node's trajectory is the pose
$\Twb^{(i)}(t)$ of $B_i$ in the node's \emph{local world frame} $W_i$,
initialized at startup. Inter-node alignment is expressed as
$T_{W_j \leftarrow W_i}$, estimated by the submap registration module
(Section~\ref{sec:registration}).

\subsection{Node Descriptor}
\label{sec:descriptor}

Each node advertises a signed, CBOR-encoded \emph{node descriptor} containing
its sensor count, intrinsics, extrinsics, compute and comm tiers, and a
public signing key. Descriptors are exchanged on join and verified before any
data is merged. This enables capability-aware merging (a survey-grade LiDAR
contributes higher-weight geometry than a toy-drone LiDAR) and prevents
unauthorized nodes from polluting the map.

\subsection{Per-Node Continuous-Time Fusion}
\label{sec:ctfusion}

The per-node front-end generalizes Coco-LIC \cite{lang2023coco} and
Gaussian-LIC2 \cite{lang2026gaussian2} to $N_c$ cameras and $N_l$ LiDARs. The
state is a non-uniform cubic B-spline over $\SEthree$ representing $\Twb^{(i)}(t)$,
with control points placed adaptively based on IMU angular and linear
acceleration magnitude, exactly as in Coco-LIC.

The sliding-window Gauss-Newton solve stacks four residual classes:

\paragraph{IMU residuals.} For each IMU measurement at time $t_k$,
\begin{equation}
  \mathbf{r}_{\text{IMU}}(t_k) =
  \begin{bmatrix}
    \mathbf{a}_k - \hat{\mathbf{a}}(t_k) \\
    \boldsymbol{\omega}_k - \hat{\boldsymbol{\omega}}(t_k)
  \end{bmatrix}_{\Sigma_{\text{IMU}}^{-1}},
\end{equation}
where $\hat{\mathbf{a}}(t_k)$ and $\hat{\boldsymbol{\omega}}(t_k)$ are derived
from the spline's analytical first and second derivatives, as in Coco-LIC.

\paragraph{LiDAR residuals.} For each LiDAR $l \in \{1,\dots,N_l\}$ and each
point $\mathbf{p}_k^{(l)}$ acquired at time $t_k$, we transform the point into
$W_i$ using the spline-evaluated pose at $t_k$ and register it against the
local Gaussian map $\GaussianMap_i$ via a point-to-Gaussian residual:
\begin{equation}
  \mathbf{r}_{\text{LiDAR}}^{(l)}(t_k) =
    \mathbf{n}^\top \bigl( \Twb^{(i)}(t_k)\, T_{B_i \leftarrow L_l}\, \mathbf{p}_k^{(l)} - \boldsymbol{\mu}_* \bigr),
\end{equation}
where $\boldsymbol{\mu}_*$ is the closest Gaussian mean in $\GaussianMap_i$
and $\mathbf{n}$ the local surface normal estimated from the neighboring
Gaussian covariances. This generalizes Gaussian-LIC2's single-LiDAR
registration to $N_l$ independent LiDAR streams sharing the same spline.

\paragraph{Camera residuals.} For each camera $c \in \{1,\dots,N_c\}$ and each
keyframe at exposure mid-time $t_k$, a photometric residual is computed by
rendering $\GaussianMap_i$ at the spline-evaluated pose and comparing against
the observed image with an SSIM+L1 combination, exactly as in Gaussian-LIC2.

\paragraph{Time-offset residuals.} Each sensor pair has an unknown time offset
$\Delta t$ estimated jointly with the spline, using Coco-LIC's online
formulation. With $N_c$ cameras and $N_l$ LiDARs there are
$N_c + N_l$ offsets relative to the IMU, all estimated in the same solve.

All residuals are weighted by their information matrices and solved in a
sliding window with the B-spline motion prior, yielding the per-node trajectory
and a locally photo-realistic Gaussian map $\GaussianMap_i$.

\subsection{Network Topology: Hierarchical Hybrid}
\label{sec:topology}

We adopt a \emph{hierarchical hybrid} topology with three logical tiers
(Figure~\ref{fig:topology}):

\begin{itemize}
  \item \textbf{Tier-0 (optional): an edge fusion server.} A fixed or
    vehicle-mounted workstation that performs global map fusion, cross-cluster
    loop-closure search, and provides a ground-truth anchor. Absent in ad-hoc
    deployments; the system degrades to Tier-1 P2P.
  \item \textbf{Tier-1: cluster heads.} Nodes self-elect cluster heads by a
    deterministic score combining compute tier and link quality. A cluster head
    merges submaps from its members, talks P2P to other heads, and uplinks to
    the Tier-0 server if present.
  \item \textbf{Tier-2: cluster members.} Plain nodes that run local
    Gaussian-LIC2 odometry and publish compressed submaps to their head.
\end{itemize}

\begin{figure}[t]
  \centering
  \resizebox{\columnwidth}{!}{% Topology figure for Co-LIC2 paper
% Hierarchical hybrid: Tier-0 (optional) edge server, Tier-1 cluster heads (P2P), Tier-2 members
\begin{tikzpicture}[
  font=\footnotesize,
  >=Stealth,
  node distance=0.6cm and 1.2cm,
  tier0/.style={draw, rounded corners, fill=blue!10, minimum width=3.2cm, minimum height=0.7cm, align=center},
  tier1/.style={draw, rounded corners, fill=green!12, minimum width=2.0cm, minimum height=0.7cm, align=center},
  tier2/.style={draw, rounded corners, fill=orange!10, minimum width=1.7cm, minimum height=0.6cm, align=center},
  bus/.style={very thick, gray},
  link/.style={->, thick},
]

% Tier-0
\node[tier0] (server) {Tier-0: Edge Fusion Server};

% Tier-1 bus
\node[draw, dashed, rounded corners, fill=gray!8, minimum width=11cm, minimum height=0.9cm,
      below=1.0cm of server] (bus1) {};
\node[anchor=west] at (bus1.west) {\scriptsize Inter-Cluster P2P Bus (DDS / gRPC)};

% Cluster heads
\node[tier1, below=0.4cm of bus1.west, xshift=1.2cm] (h1) {Cluster Head 1};
\node[tier1, below=0.4cm of bus1.center] (h2) {Cluster Head 2};
\node[tier1, below=0.4cm of bus1.east, xshift=-1.2cm] (h3) {Cluster Head 3};

% P2P links between heads
\draw[bus, <->] (h1.north) -- ($(bus1.south -| h1.north)$);
\draw[bus, <->] (h2.north) -- ($(bus1.south -| h2.north)$);
\draw[bus, <->] (h3.north) -- ($(bus1.south -| h3.north)$);
\draw[<->, thick, green!50!black] (h1.east) -- (h2.west);
\draw[<->, thick, green!50!black] (h2.east) -- (h3.west);

% Uplink from bus to server
\draw[link, blue!60!black] (server.south) -- ($(bus1.north -| server.south)$)
  node[midway, right, font=\scriptsize] {uplink (compressed)};

% Tier-2 members
\node[tier2, below=0.7cm of h1] (m1a) {Node 1};
\node[tier2, below=0.7cm of m1a] (m1b) {Node 2};
\node[tier2, below=0.7cm of h2] (m2a) {Node 3};
\node[tier2, below=0.7cm of m2a] (m2b) {Node 4};
\node[tier2, below=0.7cm of h3] (m3a) {Node 5};
\node[tier2, below=0.7cm of m3a] (m3b) {Node 6};

% member -> head links
\draw[link, orange!70!black] (m1a.north) -- (h1.south);
\draw[link, orange!70!black] (m1b.north) -- (m1a.south);
\draw[link, orange!70!black] (m2a.north) -- (h2.south);
\draw[link, orange!70!black] (m2b.north) -- (m2a.south);
\draw[link, orange!70!black] (m3a.north) -- (h3.south);
\draw[link, orange!70!black] (m3b.north) -- (m3a.south);

% Legend
\node[draw, fill=blue!10, rounded corners, minimum width=1.4cm, anchor=west]
  at ($(server.north east)+(0.3,0.5)$) (l1) {\scriptsize Tier-0};
\node[draw, fill=green!12, rounded corners, minimum width=1.4cm, anchor=west]
  at (l1.east) (l2) {\scriptsize Tier-1 head};
\node[draw, fill=orange!10, rounded corners, minimum width=1.4cm, anchor=west]
  at (l2.east) (l3) {\scriptsize Tier-2 node};

\end{tikzpicture}
}
  \caption{Hierarchical hybrid topology. An optional Tier-0 edge server fuses
  cross-cluster submaps; cluster heads (Tier-1) merge their members' submaps and
  talk P2P; cluster members (Tier-2) run local Gaussian-LIC2 and publish
  submaps. On server loss the system degrades to Tier-1 P2P; on head loss,
  members re-elect.}
  \label{fig:topology}
\end{figure}

\paragraph{Cluster-head election.} Each node computes a score
\begin{equation}
  \text{score}(i) = w_c \cdot \text{compute\_tier}(i) + w_q \cdot \sqrt{\text{link\_quality}(i)} + w_s \cdot N_s^{(i)},
\end{equation}
where $N_s^{(i)} = N_c^{(i)} + N_l^{(i)}$ is the total sensor count, and
$w_c, w_q, w_s$ are weights (defaults $w_c = 0.5, w_q = 0.4, w_s = 0.05$). The
highest-scoring node in a 2-hop neighborhood becomes head; ties are broken by
$\text{node\_id}$. Re-election triggers on heartbeat timeout (default $3$~s).

\subsection{Software Architecture}
\label{sec:software}

Each node runs a ROS~2 / DDS component graph. The local pipeline is
Gaussian-LIC2 with new cooperative components: a \emph{submap extractor} that
carves a spatial sub-volume of Gaussians, a \emph{submap compressor} with a
budget controller, a \emph{communication layer} (DDS with gRPC fallback), a
\emph{map merger} that aligns and reconciles incoming submaps, a
\emph{loop-closure agent}, and a \emph{coordinator} that handles discovery and
election. Module responsibilities and reuse from Gaussian-LIC2 are tabulated in
the companion design document.

\begin{figure}[t]
  \centering
  \resizebox{\columnwidth}{!}{% Data-flow figure for Co-LIC2 paper
% Per-node pipeline: sensors -> CT fusion -> local map -> submap pub; incoming -> merge -> PGO -> correction
\begin{tikzpicture}[
  font=\footnotesize,
  >=Stealth,
  node distance=0.5cm and 0.9cm,
  box/.style={draw, rounded corners, minimum width=2.4cm, minimum height=0.7cm, align=center},
  sens/.style={box, fill=gray!10},
  core/.style={box, fill=blue!10},
  map/.style={box, fill=green!12},
  coop/.style={box, fill=orange!12},
  pgo/.style={box, fill=purple!12},
  arr/.style={->, thick},
]

% Sensor column
\node[sens] (cam) {Cameras $N_c$};
\node[sens, below=of cam] (lid) {LiDARs $N_l$};
\node[sens, below=of lid] (imu) {IMU};

% Sync
\node[box, fill=gray!15, right=1.0cm of lid] (sync) {Sync \&\\time-stamp};

% CT fusion
\node[core, right=1.0cm of sync] (ct) {CT B-spline\\fusion (multi-sensor)};

% Trajectory + local map
\node[box, fill=yellow!15, above right=0.3cm and 1.0cm of ct] (traj) {Trajectory $\Twb(t)$};
\node[map, below right=0.3cm and 1.0cm of ct] (gmap) {Local Gaussian\\map $\GaussianMap_i$};

% Submap publish path
\node[coop, right=1.0cm of gmap] (ext) {Submap\\extractor};
\node[coop, right=1.0cm of ext] (cmp) {Submap\\compressor};
\node[coop, right=1.0cm of cmp] (comm) {Comm layer\\$\to$ head / peers};

% Incoming merge path
\node[coop, below=1.0cm of comm] (minc) {Incoming\\submaps};
\node[coop, left=1.0cm of minc] (merge) {Map merger};
\node[pgo, left=1.0cm of merge] (pgo) {Pose graph\\(iSAM2)};

% Correction back to local map
\node[box, fill=red!10, below=1.0cm of pgo] (corr) {Apply PGO\\correction};

% Edges
\draw[arr] (cam.east) -- (sync.west);
\draw[arr] (lid.east) -- (sync.west);
\draw[arr] (imu.east) -- (sync.west);
\draw[arr] (sync.east) -- (ct.west);
\draw[arr] (ct.east) -- (traj.west);
\draw[arr] (ct.east) -- (gmap.west);
\draw[arr] (traj.east) to[bend left=20] (ext.north);
\draw[arr] (gmap.east) -- (ext.west);
\draw[arr] (ext.east) -- (cmp.west);
\draw[arr] (cmp.east) -- (comm.west);

\draw[arr] (minc.west) -- (merge.east);
\draw[arr] (merge.west) -- (pgo.east);
\draw[arr] (pgo.south) -- (corr.north);
\draw[arr, dashed, red!60!black] (corr.east) to[bend right=20] (gmap.south);
\draw[arr, dashed, red!60!black] (merge.north) to[bend left=15] (gmap.south east);

% Loop closure annotation
\node[coop, below=0.4cm of merge] (loop) {Loop-closure agent};
\draw[arr, dashed] (loop.north) -- (merge.south);
\draw[arr, dashed] (loop.west) -- (pgo.south east);

\end{tikzpicture}
}
  \caption{Per-node data flow. Sensor drivers feed the CT fusion front-end,
  which produces a trajectory and updates the local Gaussian map. The submap
  extractor compresses a spatial slice and publishes it; incoming submaps are
  merged back into the local map; loop-closure constraints feed a pose graph
  that broadcasts corrections.}
  \label{fig:dataflow}
\end{figure}
